\section{第一章 绪论}\label{ux7b2cux4e00ux7ae0-ux7eeaux8bba}

\subsection{1.1
研究背景与意义}\label{ux7814ux7a76ux80ccux666fux4e0eux610fux4e49}

现场可编程门阵列（Field-Programmable Gate
Array，FPGA）是一种可在出厂后由用户重复配置其内部逻辑结构的可编程逻辑器件。由于其兼顾了专用集成电路（ASIC）的高性能与通用处理器的灵活性，FPGA
在通信基础设施、人工智能推理加速、图像处理、密码学硬件以及嵌入式控制等领域中得到了极为广泛的应用{[}1{]}{[}2{]}。特别是在深度学习加速与边缘计算领域，FPGA
凭借其低延迟、低功耗和可重配置特性，已成为替代 GPU
进行定制化推理的重要选项{[}3{]}。

然而，FPGA 的设计与验证流程本质上是一项高度人工密集的工程任务。传统的
FPGA
设计流程涵盖行为级规格描述、寄存器传输级（RTL）代码编写、功能仿真、逻辑综合、布局布线以及比特流烧写等多个阶段{[}4{]}。工程师不仅需要熟练掌握硬件描述语言（如
Verilog 或 VHDL），还需要深入理解目标器件的微架构约束、时序特性以及 EDA
工具链的使用规范。从规格到验证通过的完整开发周期，即便对于一个中等规模的模块，也往往需要数天乃至数周的人工投入。

\begin{figure}
\centering
\pandocbounded{\includegraphics[keepaspectratio,alt={图1-1: 系统高层工作流}]{../fpga_flow/diagrams/01_high_level_workflow.png}}
\caption{图1-1: 系统高层工作流}
\end{figure}

2022 年以来，以 GPT-4、Claude 以及 DeepSeek 为代表的大语言模型（Large
Language
Model，LLM）在代码生成领域取得了突破性进展{[}5{]}{[}6{]}。这些模型不仅能够生成逻辑正确的高级语言代码，也逐步展现出生成硬件描述语言的能力。与此同时，多智能体系统（Multi-Agent
System，MAS）理论与工程实践的快速发展，使得将多个专业化 LLM Agent
组织为协作流水线、分工完成复杂工程任务成为可能{[}7{]}。AutoGen{[}8{]}、CrewAI、OpenHands
SDK{[}9{]}等框架的涌现，为构建以 LLM
为核心的自动化设计流程提供了基础设施支撑。

将 LLM 与多智能体协作机制相结合，应用于 FPGA
设计自动化，具有重要的研究价值与现实意义：一方面，自动化的代码生成与验证闭环可以大幅降低
RTL 设计的人工成本；另一方面，通过系统化的约束工程与验证层级设计，可以将
LLM
的不确定性和幻觉风险约束在可控范围内，从而实现从规格输入到验证通过的全流程自动化。本文设计并实现了一套通用的多智能体
FPGA 设计全流程自动化框架，其核心架构（编排状态机、分层验证体系、LLM
约束引擎）适用于任何可分解为 DAG 结构的多模块硬件设计任务，并以 AES-128
加密硬件全系统作为验证案例，端到端地验证了该框架的有效性与工程可行性。

\subsection{1.2
国内外研究现状}\label{ux56fdux5185ux5916ux7814ux7a76ux73b0ux72b6}

\subsubsection{1.2.1 LLM
用于硬件描述语言生成}\label{llm-ux7528ux4e8eux786cux4ef6ux63cfux8ff0ux8bedux8a00ux751fux6210}

近年来，学界已有多项工作探索 LLM 在 RTL
代码生成方面的能力边界。VerilogEval（Liu
等，2023）{[}10{]}建立了系统性的 Verilog 代码生成评测基准，发现 GPT-4
在功能正确性上显著优于早期模型，但在复杂的时序逻辑和多模块协同设计方面仍存在明显局限。RTLLM（Lu
等，2024）{[}11{]}进一步扩展了评测范围，引入了综合后的时序收敛指标，揭示了单纯以仿真通过为标准可能掩盖的设计质量问题。ChipChat（Blocklove
等，2023）{[}12{]}则展示了通过人机交互对话生成 RISC-V
处理器的可行性，但其过程需要大量人工干预，尚未实现真正意义上的自动化。

上述工作共同揭示了一个基本规律：LLM 在单模块、功能相对简单的 RTL
生成任务中表现尚可，但在面对多模块相互依赖的系统级设计时，由于缺乏跨模块的一致性约束和自动化的验证闭环，生成结果的可靠性会显著下降。

\subsubsection{1.2.2 多智能体框架与 EDA
自动化}\label{ux591aux667aux80fdux4f53ux6846ux67b6ux4e0e-eda-ux81eaux52a8ux5316}

在多智能体框架方面，AutoGen（Wu 等，2023）{[}8{]}提出了基于可对话 Agent
的多智能体协作范式，通过 ConversableAgent
链式调用实现了多步骤任务的自动化执行。OpenHands（Wang
等，2024）{[}9{]}则提供了一套以 Agent-Conversation-Tool 为核心抽象的
SDK，支持构建具有工具调用能力的工作流编排系统。

在 EDA 智能化方向，ChipGPT 探索了以 LLM 为核心驱动 EDA
工具链的可能性，ChatEDA 尝试以自然语言接口封装 Synopsys
工具链，DAMO-Chip 则关注流程中的知识辅助与错误诊断。这些工作均验证了 LLM
与 EDA
工具集成的技术路径，但大多停留在单工具调用或单阶段自动化层面，缺乏对完整设计验证流程的系统性覆盖{[}13{]}。

\subsubsection{1.2.3
同校前序工作}\label{ux540cux6821ux524dux5e8fux5de5ux4f5c}

本文的直接参考对标是同校同专业的前序毕业设计工作{[}14{]}，该工作基于
AutoGen 框架，构建了一条以 AutoGen ConversableAgent
为核心的线性六步流水线，依次完成需求解析、方案设计、RTL 代码生成、TCL
脚本生成、Vivado 仿真验证以及结果评估等阶段，并引入了 DeepSeek-reasoner
与本地微调代码生成模型（codegen-2B）协同工作的双模型策略。该工作证明了多
Agent 协作路径用于 FPGA
自动化设计的基本可行性，为本文的研究提供了重要参考基础。

然而，该工作在架构层面存在若干系统性局限，这些局限在面对更复杂的多模块协同设计任务时会成为制约瓶颈，本文第
1.3 节将对此进行详细分析。

\subsection{1.3
现有方案的不足}\label{ux73b0ux6709ux65b9ux6848ux7684ux4e0dux8db3}

通过对基于AutoGen的线性pipeline方案以及上述同类工作的系统性分析，本文归纳出现有方案在以下五个维度上存在明显不足：

\textbf{不足一：Prompt-controlled 方案的可靠性问题。}
AutoGen方案以及大多数同类工作均属于 Prompt-controlled
架构------系统的执行进度、分支决策和终止条件均由 LLM
的提示词输出决定，框架本身不拥有控制权。这意味着 LLM
的幻觉行为（hallucination）可以直接影响系统执行路径：LLM
可能生成不存在的端口名称、错误的模块接口，或在验证未完成时提前宣布任务成功。在没有外部强制验证机制的情况下，这类错误难以被系统自动检测和拦截。

\textbf{不足二：线性 pipeline 无法处理模块间依赖。}
AutoGen方案采用的六步线性 pipeline
假设设计任务可以被分解为一条无分支的顺序执行链。然而，实际的多模块 FPGA
设计往往具有复杂的模块间依赖关系。以本文的 AES-128
案例为例，\texttt{aes\_sbox}
模块被另外三个模块所依赖，\texttt{aes\_round\_transform} 与
\texttt{aes\_key\_schedule\_128} 共同服务于顶层的
\texttt{aes128\_encrypt\_core}。线性 pipeline
无法表达和调度这种有向无环图（DAG）结构，既无法利用拓扑并行性加速执行，也无法正确处理依赖前置与批次隔离的需求。

\textbf{不足三：缺乏结构化修复机制。} 基于AutoGen的前序方案中的 Debug
Agent
采用''读取仿真日志→重写代码''的模式处理验证失败。这种方式存在两个根本性问题：其一，修复行为缺乏约束------Agent
可以对代码进行任意范围的修改，容易引入新的错误；其二，修复循环缺乏收敛保证------在没有预算限制和升级策略的情况下，系统可能陷入无限修复循环而无法终止。

\textbf{不足四：验证层级单一。} AutoGen方案仅依赖 Vivado xsim
仿真作为唯一的验证手段，通过波形图人工判读结果。这种单层验证策略有两个局限：首先，语法错误只有在完整仿真时才被发现，缺少编译门控层的快速失败机制；其次，仿真通过仅验证了已知的测试向量，缺乏对随机输入和边界条件的鲁棒性检验，也没有针对多模块协同场景下的集成回归测试。

\textbf{不足五：对 LLM 的系统性约束工程不足。}
AutoGen方案和大多数同类工作对 LLM 的约束主要依赖
\texttt{system\_message} 中的自然语言指令，本质上是对 LLM
的''软性请求''。当 LLM 在复杂场景下偏离指令（scope
drift、格式错误、越界修改）时，系统缺乏机器可检测的约束机制加以拦截。在工程可靠性要求较高的硬件设计场景中，这种''仅靠
Prompt 约束''的策略存在根本性的可靠性隐患。

\subsection{1.4
本文主要工作与创新点}\label{ux672cux6587ux4e3bux8981ux5de5ux4f5cux4e0eux521bux65b0ux70b9}

针对上述不足，本文设计并实现了一套通用的多智能体 FPGA
设计全流程自动化框架，提出以下三个核心创新点。三个创新点均以通用硬件设计自动化为目标进行架构设计，以
AES-128 加密硬件全系统作为验证案例进行端到端实验验证：

\textbf{创新点一：框架控制的编排状态机与 AgentExecutionPolicy
策略分层架构。}
本文提出一种''框架控制''（Framework-controlled）的多智能体编排模式，通过参数化的编排状态机精确定义硬件设计自动化的通用执行阶段（规格摄入→架构设计→规划→模块设计→分层验证→集成回归→终态）与转移条件，将状态推进权从
LLM 手中转移到框架本身。配合 \texttt{AgentExecutionPolicy}
策略层，将每个编排状态映射到特定的 LLM 配置（thinking/fast 双
Profile）和子 Agent
委派策略（forbidden/module\_worker\_allowed/l2\_campaign\_allowed），实现了对
Agent
行为边界的精确控制。同时设计了基于修复预算的升级策略，当模块工作者失败超过阈值时自动升级至更高推理能力的编排器，保证了修复循环的收敛性。在
AES-128 验证案例中，该状态机实例化为 10 个具体状态，覆盖 4 模块 DAG
的完整生命周期。

\textbf{创新点二：L0→L1→L2→Integration 四层分级验证体系与 Checkpoint
协议。}
本文设计了一套适用于任何多模块硬件设计的由浅入深的四层分级验证架构：L0
编译门控层负责秒级语法检查，L1 仿真与 Checkpoint
层实现已知向量的功能验证，L2 鲁棒性测试层覆盖随机和边界用例，Integration
集成回归层验证多模块协同行为。各层之间形成严格的门控关系，前一层验证通过才能进入下一层。Checkpoint
协议采用
\texttt{CHECKPOINT\textbar{}\textless{}name\textgreater{}\textbar{}PASS\textbar{}\textless{}detail\textgreater{}}
的结构化文本格式，由 testbench
主动发射、框架自动解析，将验证结论的判定权从人工判读转移到机器自动处理。该验证体系的架构设计不依赖于特定硬件设计------任何具备
Verilator 可编译 RTL 和自检 Testbench 的设计均可接入。在 AES-128
验证案例中，4 个模块共定义了 9 个 Checkpoint，集成回归测试覆盖了
baseline、back\_to\_back、mid\_reset 三种 Campaign。

\textbf{创新点三：Harnessing Engine --- 六层 LLM 约束引擎。} 本文提出
Harnessing Engine（驾驭引擎）概念，作为一种通用的 LLM
行为约束框架，以结构化工程手段从外部约束 LLM
在硬件设计任务中的行为，而非依赖 LLM
的''自律性''。该引擎由六个层级构成：Skill 层注入领域知识，Prompt
层通过可组合合约规范 Agent 行为边界，Memory
层将验证通过的参考代码预填充到工作区以降低生成难度，Hooks
层通过自定义工具代理所有执行操作并实施哈希验证，MCP 层通过 Model Context
Protocol 标准化 Agent 与 EDA 工具的通信并过滤危险工具，SlashCommand
层提供语义完整的 CLI
命令体系。六层架构本身是领域无关的约束模式，可应用于任何需要驾驭 LLM
不确定性的工程场景。在 AES-128 验证案例中，Skill 层实例化为 16 个 AES
领域技能文件，Memory 层预填充 4 个模块的验证通过参考 RTL 与
Testbench，针对 LLM
在硬件设计中的典型失败模式（幻觉端口、语法错误、scope
drift、过度修改、格式错误）构建了多重防御屏障。

以上三个创新点在架构层面均以通用多模块硬件设计自动化为目标设计，不依赖于特定的设计对象。本文以
AES-128 加密全系统（包含
\texttt{aes\_sbox}、\texttt{aes\_key\_schedule\_128}、\texttt{aes\_round\_transform}、\texttt{aes128\_encrypt\_core}
四个模块组成的
DAG）作为主要验证案例，端到端地实验证明了该通用框架的有效性，实现了从规格输入到全部
Checkpoint 通过的全流程自动化。此外，本文以 LED Chaser（4 位 LED
流水灯）作为第二个验证案例，通过 YAML
蓝图驱动的方式在不修改任何框架代码的前提下完成了从规格输入到集成验证通过的全流程执行，进一步实证了框架的跨设计通用性。

\subsection{1.5
论文组织结构}\label{ux8bbaux6587ux7ec4ux7ec7ux7ed3ux6784}

本文共分八章，各章内容安排如下：

第二章介绍本文所涉及的相关技术与理论基础，包括大语言模型与代码生成、多智能体系统理论、FPGA
设计与验证基础，以及 AES-128 加密算法原理。

第三章对框架进行总体设计概述，从通用架构设计的宏观视角描述六层框架结构、AES-128
模块 DAG 分解（作为框架实例化的验证案例）、Agent 角色与 LLM
配置策略，以及 DAG 拓扑批调度机制。

第四章深入展开编排状态机与策略分层设计，详细分析 10
状态状态机的定义与转移逻辑、AgentExecutionPolicy
的策略映射、升级机制设计，以及框架控制与 Agent 控制的根本差异。

第五章深入展开通用四层验证体系与 Checkpoint
协议，分析各验证层级的设计动机与通用性边界、Checkpoint
协议的工作机制、L2 鲁棒性测试策略以及集成回归测试的三种 Campaign。

第六章深入展开 Harnessing Engine 的通用 LLM 约束框架设计，逐层分析
Skill、Prompt、Memory、Hooks、MCP 和 SlashCommand
各层的约束机制与协同效果，并说明每层约束的领域无关性。

第七章以 AES-128 和 LED Chaser 双案例呈现框架的实验验证结果，包括
AES-128 单模块验证数据、集成回归测试结果、全流程执行分析、与 AutoGen
方案的能力维度对比，以及 LED Chaser 跨设计通用性验证。

第八章总结本文通用框架设计与三个创新点的主要贡献，并结合 AES-128
验证案例的实验结果指出当前局限与未来研究方向。
